% Generated from locale/tr/content/turing-machines/undecidability/halting-problem.tex; do not hand-edit.


\olfileid[tr]{tur}{und}{hal}
\olsection{Durma Problemi}

\begin{explain}
Turing makinesi betimlerinin bir numaralandırmasını sabitlediğimizi
varsayalım. Böylece her Turing makinesi bir \emph{indis}, yani Turing
makinesi betimlerinin $M_1$, $M_2$, $M_3$, \dots{} numaralandırmasındaki
yerini alır.

Turing-hesaplanabilir olmayan fonksiyonların bulunması gerektiğini biliyoruz:
Turing makinesi betimleri kümesi ve dolayısıyla Turing makineleri kümesi
!!{enumerable}, ama $\Nat$'ten $\Nat$'e bütün fonksiyonlar kümesi değildir.
Bununla birlikte hesaplanamaz fonksiyonların belirli örneklerini de
bulabiliriz. Bunlardan biri durma fonksiyonudur.
\end{explain}

\begin{defn}[Durma fonksiyonu]
 \emph{Durma fonksiyonu}~$h$ şöyle tanımlanır:
\[
h(e,n) =
\begin{cases}
  \text{0} & \text{$M_e$ makinesi $n$ girdisinde durmazsa} \\
  \text{1} & \text{$M_e$ makinesi $n$ girdisinde durursa}
\end{cases}
\]
\end{defn}

\begin{defn}[Durma problemi]
\emph{Durma Problemi}, her $e$, $n$ için $M_e$ Turing makinesinin $n$ çizgili
bir girdide durup durmadığını belirleme problemidir.
\end{defn}

\begin{explain}
İlişkili bir $s$ fonksiyonunun Turing-hesaplanabilir olmadığını göstererek
$h$'nin de Turing-hesaplanabilir olmadığını göstereceğiz. Bu kanıt, bir
Turing makinesince hesaplanabilen her şeyin disiplinli bir Turing makinesince
hesaplanabileceği (\olref[mac][dis]{sec}) ve iki Turing makinesinin tek bir
makine oluşturacak biçimde birbirine bağlanabileceği
(\olref[mac][cmb]{sec}) olgularına dayanır.
\end{explain}

\begin{defn} $s$ fonksiyonu şöyle tanımlanır:
\[
s(e) =
\begin{cases}
  \text{0} & \text{$M_e$ makinesi $e$ girdisinde durmazsa} \\
  \text{1} & \text{$M_e$ makinesi $e$ girdisinde durursa}
\end{cases}
\]
\end{defn}

\begin{lem}
$s$ fonksiyonu Turing hesaplanabilir değildir.
\end{lem}

\begin{proof}
Çelişki elde etmek üzere $s$ fonksiyonunun Turing hesaplanabilir olduğunu
varsayalım. O zaman $s$'yi hesaplayan bir $S$ Turing makinesi bulunurdu.
Genelliği yitirmeden $S$ durduğunda bunu ilk kareyi tararken yaptığını (yani
disiplinli olduğunu) varsayabiliriz. Bu makine, $0$ girdisinde
başlatıldığında duran (yani ilk durumda şerit sonu simgesinin sağındaki kareyi
tararken $\TMblank$ okursa duran), aksi durumda sağa doğru uzaklaşıp hiç
durmayan başka bir $J$ makinesine ``bağlanabilir''. $S$ ile $J$'nin
bağlanmasıyla oluşturulan $S\concat J$ bir Turing makinesidir; dolayısıyla
bir $e$ için $M_e$'dir (yani numaralandırmanın bir yerinde bulunur). $M_e$'yi
$e$ tane $\TMstroke$ girdisinde başlatın. İki olasılık vardır: $M_e$ ya durur
ya da durmaz.
\begin{enumerate}
\item $M_e$'nin $e$ tane $\TMstroke$ girdisinde durduğunu varsayalım. O zaman
  $s(e)=1$'dir. Dolayısıyla $S$, $e$ girdisinde başlatıldığında şeritte çıktı
  olarak tek bir $\TMstroke$ ile durur. Ardından $J$, şeritte bir
  $\TMstroke$ ile başlar. Bu durumda $J$ durmaz. Ancak $M_e$, $S\concat J$
  makinesidir; dolayısıyla önce $S$'nin, sonra $J$'nin yapacağını yapmalıdır
  (yani bu durumda sağa doğru uzaklaşıp hiç durmamalıdır). Öyleyse $M_e$,
  $e$ tane $\TMstroke$ girdisinde duramaz.

\item Şimdi $M_e$'nin $e$ tane $\TMstroke$ girdisinde durmadığını
  varsayalım. O zaman $s(e)=0$'dır ve $S$, $e$ girdisinde başlatıldığında boş
  bir şeritle durur. $J$ boş bir şeritte başlatıldığında hemen durur. Yine
  $M_e$, önce $S$'nin sonra $J$'nin yapacağını yapar; dolayısıyla $M_e$, $e$
  tane $\TMstroke$ girdisinde durmalıdır.
\end{enumerate}
Her iki durumda da varsayımımızla çelişkiye ulaşırız. Bu, böyle bir $S$
Turing makinesinin bulunamayacağını gösterir: $s$ Turing hesaplanabilir
değildir.
\end{proof}

\begin{thm}[Durma Probleminin Çözülemezliği]
\ollabel{thm:halting-problem} Durma problemi çözülemezdir; yani $h$ fonksiyonu
Turing hesaplanabilir değildir.
\end{thm}

\begin{proof}
$h$'nin, diyelim bir $H$ Turing makinesince, Turing hesaplanabilir olduğunu
varsayalım. $H$'yi $s$'yi hesaplayan bir Turing makinesi kurmak için
kullanabilirdik: Önce girdinin bir kopyasını çıkarın (kopyaları bir
$\TMblank$ simgesiyle ayırın). Sonra başa dönüp $H$'yi çalıştırın. İlk işlemi
yapan bir makineyi açıkça kurabiliriz (bkz.~\cref{tur:mac:dis:prob:copier});
$H$ var olsaydı böyle bir kopyalama makinesine onu ``bağlayarak'' $M_e$'nin
$e$ girdisinde durup durmadığını belirleyen, yani $s$'yi hesaplayan yeni bir
makine elde edebilirdik. Fakat böyle bir makinenin bulunamayacağını zaten
gösterdik. Öyleyse $h$ de Turing hesaplanabilir değildir.
\end{proof}

\begin{prob}
Üçte Durma (3-Durma) problemi, keyfî seçilmiş bir Turing makinesinin bunun
dışında boş olan bir şeritte üç $\TMstroke$ girdisi için durup durmadığını
belirleyen bir karar süreci verme problemidir. 3-Durma probleminin çözülemez
olduğunu kanıtlayın.
\end{prob}

\begin{prob}
Durma problemi, $e\neq n$ olan Turing makinesi ve girdi çiftleri $M_e$ ile
$n$ için çözülebilirse $e=n$ durumları için de çözülebileceğini gösterin.
\end{prob}

\begin{prob}
Girdinin, bütün Turing makinelerinin bir numaralandırması aracılığıyla bir
$M_e$ Turing makinesini belirleyen $e$ sayısı olması durumunda durma
probleminin çözülemez olduğunu kanıtladık. Ya
\olref[tur][und][enu]{sec} içindeki Turing makinesi betimlerini doğrudan
girdi olarak kullanmaya izin verirsek? İndisler yerine Turing makinelerinin
betimlerini girdi olarak alan ve durma problemine karar veren bir Turing
makinesi bulunabilir mi? Nedenini açıklayın.
\end{prob}

\begin{prob} Aşağıdaki gibi tanımlanan \emph{kısmi} $s'$ fonksiyonunun
  \[
  s'(e) =
  \begin{cases}
    \text{1} & \text{$M_e$ makinesi $e$ girdisinde durursa}\\
    \text{tanımsız} & \text{$M_e$ makinesi $e$ girdisinde durmazsa}
  \end{cases}
  \]
  Turing hesaplanabilir \emph{olduğunu} gösterin.
\end{prob}

